\section{Introduction}
\label{sec:1}

\subsection{Motivation}
\label{sec:1.1}

Modern large-scale distributed software systems progressively use asynchronous, event-driven, publish--subscribe (pub-sub) architectures. These are common within autonomous driving (ROS~2 \cite{macenski2022robot}), enterprise event streams (Apache Kafka \cite{kreps2011kafka}), cyber-physical systems (DDS \cite{omg2015dds}), IoT deployments (MQTT \cite{oasis2019mqtt}), cloud-native microservices \cite{dragoni2017microservices,newman2015building}, and distributed AI/LLM clusters. Pub-sub architectures separate producers and consumers in terms of space, time, and synchronization \cite{eugster2003faces}. Components communicate indirectly through message topics and brokers, so they do not need direct static references. Current middleware also supports deployment-time Quality-of-Service (QoS) policies, such as reliability, durability, message priorities, and delivery deadlines, to help manage performance during peak loads and network stress.

While this separation allows for flexible scaling, it also makes it harder to see what is happening inside the system. Unlike synchronous architectures like RESTful HTTP or gRPC, which show clear caller and callee paths, pub-sub publishers and subscribers lack direct links. As a result, failures, head-of-line blocking, and backpressure can spread along concealed paths involving brokers, shared topics, colocated hosts, and shared libraries \cite{motter2002cascade,buldyrev2010cascade}. These failures mainly happen in two ways: sequential cascades, where a slow subscriber fills up a broker queue and gradually slows down its publishers \cite{albert2000error}; and simultaneous blast radii, where a shared library crash or host outage brings down all colocated services. Standard architecture diagrams and static call graphs cannot show these processes. The best time to reduce these risks is before deployment, especially during design and continuous integration, following dependable computing principles \cite{avizienis2004basic,bass2012software}. At these points, runtime telemetry, distributed tracing, and logs are not available. Therefore, architects and Site Reliability Engineers need to identify the most critical components, topics, and links, understand why they matter, and take targeted actions like broker replication, splitting overloaded topics, or isolating shared libraries to reduce risk. Here, systemic criticality means an entity is either a single point of failure that can disrupt downstream connections if it fails, or an error-propagation hub that can cause large cascading outages in the asynchronous message mesh.

By analyzing architecture directly from configuration manifests, pre-deployment Static System Analysis removes the need for staging clusters, active containers, and live fault-injection tools, which reduces infrastructure overhead. However, static analysis is not faster than in-process simulation in terms of CPU time: extracting topology features takes $79.3\,\text{s}$ for a 520-component mesh and $239.3\,\text{s}$ for 2,000 components. In contrast, the in-process discrete-event simulator finishes in $0.08$--$4.5\,\text{s}$ (Sections~\ref{sec:7.5.oracle} and~\ref{sec:8.2}). The learned GNN forward pass is quick ($56\,\text{ms}$), but static analysis is $2$--$18$ times slower than simulation across twelve scenarios (median $5.6$ times, with the largest mesh being the slowest). This shows that static analysis mainly helps by avoiding infrastructure setup, not by being faster, unless it uses incremental topology caching. So why use static graph learning instead of simulating manifests? As explained below, manifest-level simulation needs execution details like publication rates, payload sizes, and queue capacities, which are often missing or uncalibrated during early design and integration. Message simulation also cannot directly assess unrunnable physical hosts or network links without complicated failure models. Static graph learning, by contrast, captures structural patterns and generalizes across both runnable and unrunnable infrastructure.

\subsection{Problem Statement: The Architecture--Code Gap and the Black-Box AI Challenge}
\label{sec:1.2}

Pre-deployment dependability analysis in distributed architectures has two main tasks. The first is predictive forecasting: whether a data-driven, relation-specific graph neural network can predict cascading failure blast radii and identify critical components better than traditional topological metrics. While closed-form metrics show global network connectivity, it remains unclear whether they can track multi-hop, relation-dependent cascade spread across several channels. We test this directly against such a baseline (Section~\ref{sec:7.1}). We train and evaluate the predictive model using independent simulation ground truth to rank and identify critical sets.

Explainable Criticality Attribution (the Explanation Layer) addresses the limits of ranking alone. A ranked list shows where risk is highest, but not how to fix it. To solve this, the predictor is combined with an interpretable structural quality profile based on ISO/IEC 25010 \cite{iso25010_2023} and ISO/IEC 25019 \cite{iso25019_2023}. This layer identifies structural risk types using standard software quality models. For example, it can distinguish between an unreplicated single point of failure, an error-propagating cascade hub, and a maintainability bottleneck. This helps guide targeted architectural fixes. The layer is used only for attribution, not for ranking.

This separation is built into the architecture, not just the presentation: both pathways use the same graph but do not share parameters, and neither is trained on the other's output. The term that could connect them is off by default and reported only as an ablation (Section~\ref{sec:4.2}). This independence lets Software-as-a-Graph (SaG) find components that are structurally central but have low operational impact, supporting a more detailed diagnosis than either pathway could provide alone.

The distance between an architecture as designed and as realized is long-established: Perry and Wolf \cite{perry1992foundations} formulated architectural erosion three decades ago, and the architectural-technical-debt literature has tracked it since \cite{cunningham1992wycash,garcia2009identifying}. What we term the \textbf{Architecture--Code Gap} specializes this concept to asynchronous middleware, where the challenge is that failure semantics were never expressible in artifacts a build pipeline can inspect: \emph{a distributed system can have pristine, bug-free source code within each service, yet remain fragile to ruinous global outages caused by hidden single points of failure (SPOFs) or mismatched middleware Quality-of-Service (QoS) contracts.} This fragility is acute in pub-sub architectures where producers and consumers interact indirectly without static references. Classical architecture evaluations (e.g., ATAM \cite{kazman1998atam,bass2012software}) rely on manual elicitation; static code analysis \cite{sonarsource2024clean,mccabe1976complexity,chidamber1994metrics,fenton2014metrics} inspects individual services in isolation; chaos engineering \cite{basiri2016chaos} requires provisioned staging clusters; and homogeneous centrality \cite{freeman1977centrality,brin1998pagerank,brandes2001faster,newman2010networks} flattens typed topologies into untyped graphs. Section~\ref{sec:2} develops each paradigm in detail.

Although machine learning has made big advances in software engineering, current AI methods for system dependability often work as black boxes. Deep neural models usually give risk scores or hidden representations lacking clear, actionable reasons for their predictions. In mission-critical software engineering, this shortage of transparency is not enough. Developers and Site Reliability Engineers need clear explanations of component vulnerabilities and architectural problems so they can refactor code or adjust infrastructure effectively.

\subsection{The Software-as-a-Graph (SaG) Approach}
\label{sec:1.3}

To address both the Architecture--Code Gap and the black-box AI problem, this work describes \textbf{Software-as-a-Graph (SaG)}, an AI-based pre-deployment \textbf{Static System Analysis (SSA)} framework for asynchronous and event-driven distributed systems. SaG provides the tools and experiments needed to determine whether static representations before deployment can identify systemically critical components and explain their systemic root causes. SaG uses a four-stage pipeline: (1)~modeling the architecture as a typed, directed multigraph over Applications, Brokers, Topics, Execution Hosts, and Shared Libraries (Section~\ref{sec:3.1}); (2)~mapping physical connections into a QoS-weighted semantic \texttt{DEPENDS\_ON} layer that captures sequential cascades and simultaneous failures (Section~\ref{sec:3.2}); (3)~training a Heterogeneous Graph Transformer (HGT) to predict multi-hop cascading blast radii (Section~\ref{sec:4}); and (4)~evaluating an explainable Reliability--Maintainability (RM) profile (Section~\ref{sec:5}) to show why components are fragile.

Importantly, SaG enforces a strict \textbf{input--label independence guarantee}: the learned models and attribution baselines use only the analytical graph $G_{\text{analysis}}$. Meanwhile, ground-truth failure impacts come from independent discrete-event simulators running on the raw structural topology $G_{\text{structural}}$ (Section~\ref{sec:4.4}). This separation keeps feature construction and evaluation ground-truth structurally separate, guaranteeing the method is valid. Section~\ref{sec:3} explains the full architectural flow and pipeline interactions (see Figure~\ref{fig:1}).

\paragraph{The Core Claim and Boundary Conditions}
The main claim of this paper is that pre-deployment Architecture-as-Code manifests contain enough relational and QoS information to predict multi-hop cascading failure blast radii and identify critical components, without running code, collecting runtime telemetry, or setting up staging clusters. In particular, relation-specific heterogeneous graph attention (\texttt{HGT-QoS}) provides the inductive bias needed to capture indirect cascade spread across message brokers, topics, and shared libraries, even when the distribution changes. However, this claim depends on two empirical boundary conditions:
\begin{enumerate}\tightlist
    \item \textbf{The Typing--QoS Substitution Effect:} Relational typing and continuous QoS edge encodings act as empirical substitutes rather than additive complements. While each independently produces substantial gains over an unweighted baseline ($+0.134$ and $+0.187$, Holm $p = 0.0015$), combining both exhibits diminishing returns because both channels supply redundant indicators of edge relation identity (Section~\ref{sec:7.2}).
    \item \textbf{Communication Paradigm Alignment:} Relational inductive bias transfers zero-shot to asynchronous publish--subscribe event meshes (where failure propagation flows forward via message starvation), but reverses on synchronous RPC call trees (where failures cascade backward via timeouts and retries), establishing that graph directional biases must correspond to architectural communication semantics (Section~\ref{sec:7.4}).
\end{enumerate}

\paragraph{Operational Guidance and Method Portfolio}
To help with effective adoption, SaG offers a set of methods customized to different software engineering situations:
\begin{itemize}\tightlist
    \item \textbf{HGT-QoS (Heterogeneous Learning):} Recommended for complex asynchronous publish--subscribe event meshes (e.g., ROS~2, DDS, MQTT) where cascading failures propagate through multi-hop broker topologies, shared libraries, and hardware hosts under unseen architecture topologies.
    \item \textbf{Topo-QoS (Closed-Form Centrality):} Recommended for lightweight, zero-training CI/CD gates that require instantaneous deterministic scoring without GPU resources or model training cost ($\rho = 0.553$).
    \item \textbf{Dual-Engine Consensus Protocol:} Recommended for production pre-deployment gating (Section~\ref{sec:8.1}). Because label-free uncertainty cannot reliably predict when an ML model struggles out-of-distribution ($\rho_s = -0.126$), SaG evaluates \texttt{HGT-QoS} and \texttt{Topo-QoS} concurrently: components in the unanimous top-$K$ intersection receive high-confidence triage, while ranking divergences trigger architectural review.
    \item \textbf{Explanation Layer ($Q(v)$ - ISO/IEC 25010):} Applied post-ranking as a diagnostic attribution instrument rather than a continuous ranker (Section~\ref{sec:5}), decomposing flagged components into standardized Reliability (Availability vs. Fault Tolerance) and Maintainability profiles to guide targeted refactoring (e.g., broker replication vs. circuit breakers).
\end{itemize}

\paragraph{Rationale for Graph Learning vs. Direct Simulation}
Since discrete-event simulation $I^*(v)$ gives the ground-truth criticality as well as finishes in $0.08$--$4.5\,\text{s}$, it is important to explain why we train a graph model. Manifest analysis avoids the need for staging clusters and chaos testing tools---this is about avoiding extra infrastructure, not saving energy (Section~\ref{sec:8.2} covers the energy used by the analysis itself). There are three main technical reasons to use graph learning instead of direct manifest simulation:
\begin{enumerate}\tightlist
    \item \textbf{Scoring Unsimulated Infrastructure Entities:} While discrete-event simulation sweeps message-passing components (Applications, Brokers, Libraries) across topics, it cannot directly evaluate unrunnable physical Execution Hosts or hardware interconnects without intricate physical failure models. Message passing over heterogeneous multigraphs naturally propagates relational representations across both simulated and unsimulated entities.
    \item \textbf{Robustness to Operational Parameter Drift:} Direct simulation depends on declared publication rates, payload sizes, and queue limits, which are often uncalibrated or only estimates during early design. Graph learning, however, captures topology-based and QoS patterns that extend across different architectures, even when operational details are not fully set.
    \item \textbf{Sub-Second Continuous Integration with Incremental Caching:} In CI/CD, using incremental caching limits feature extraction to the $k$-hop neighborhood around pull-request changes, allowing the $56\,\text{ms}$ forward pass for almost instant pre-commit checks. Without caching, static analysis is still slower than direct simulation in all tested scenarios.
\end{enumerate}

Section~\ref{sec:7.1} evaluates whether graph learning offers ranking advantages over closed-form baselines.

\subsection{Research Questions}
\label{sec:1.4}

This empirical study considers five research questions:

\begin{itemize}\tightlist
\item \textbf{RQ1 (Predictive Efficacy):} \emph{How accurately does heterogeneous graph learning predict cascading failure impact and identify the critical component set, compared with traditional, non-learning network indicators?}
\item \textbf{RQ2 (Value of Architectural Typing):} \emph{Does modeling distinct entity and dependency types yield better failure predictions than homogeneous graph models on architectures the model has never seen --- and does whatever advantage it confers depend on what other relational signal the model already has?}
\item \textbf{RQ3 (QoS Encoding and Robustness):} \emph{(i)~Do middleware Quality-of-Service contracts carry signal a purely structural score discards, (ii)~does that signal compose with or substitute for architectural typing, (iii)~do the framework's simulation oracles agree with one another, and (iv)~are the reported orderings robust to the free parameters of the scorer and of the ground truth?}
\item \textbf{RQ4 (Real-World Generalization):} \emph{How effectively does the framework transfer zero-shot to authentic, real-world distributed systems across autonomous driving (ROS~2), cloud-native microservices, smart home IoT, and industrial edge computing?}
\item \textbf{RQ5 (Analysis Cost):} \emph{What does pre-deployment analysis cost at CI/CD time, which pipeline stage dominates that footprint, and how does it compare against the discrete-event simulation it is intended to displace?}
\end{itemize}

\subsection{Key Contributions}
\label{sec:1.5}

This paper makes three main contributions:

\begin{enumerate}\tightlist
    \item \textbf{Heterogeneous Graph Learning for Dependability, and Its Boundaries:} A relation-specific Heterogeneous Graph Transformer that forecasts cascading blast radii from Architecture-as-Code manifests, incorporating 16-D QoS edge encodings and an auxiliary component reliability head (Section~\ref{sec:4}). Ablated separately under inductive distribution shift across twelve architectures, relation-specific HGT demonstrates out-of-distribution gains ($\Delta\rho = +0.234$ over an untyped baseline, $p = 0.0005$, reflecting the joint transition to the parameterized bidirectional HGT architecture). Crucially, typing and continuous QoS encodings act as empirical substitutes rather than complements --- an interaction that survives transformation to a variance-stabilized scale, so it is a property of the mechanisms rather than of bounded $\rho$ --- though the typing factor also carries a $15.4\times$ capacity gap and a change in message-passing directionality that the reported controls do not separate. Against an unparameterized QoS-weighted centrality baseline, learned ranking does not establish a statistically significant advantage ($+0.085$, $p = 0.151$). That parity should be read against the target: $I^*(v)$ is a deterministic functional of the same graph the predictors read, so a closed-form centrality is an expected-strong comparator. The result bounds this design more than it bounds graph learning. We report the substitution effect and the empirical boundary as primary findings (Sections~\ref{sec:7.1}--\ref{sec:7.2}).

    \item \textbf{A Formal Typed Architecture Model:} A multigraph representation that derives logical dependencies from physical pub-sub linkages and distinguishes sequential cascade propagation from simultaneous multi-consumer library failures (Section~\ref{sec:3}).

    \item \textbf{Empirical Benchmark, Real-World Transfer, and Cost Profile:} A study across seventeen system architectures totaling 2,812 components: twelve synthetic topologies (2,461 components across LOSO folds) and five open-source reference systems (351 components) under strict graph-view separation. We characterize pipeline costs and show that the neural forward pass accounts for only $0.02\%$ of runtime. In contrast, cold deterministic feature extraction dominates and costs $2$--$18\times$ the simulation it was meant to displace (median $5.6\times$ over twelve scenarios, $79.3\,\text{s}$ against $4.5\,\text{s}$ at its maximum), refuting the premise that static analysis is inherently faster than in-process simulation without incremental caching (Sections~\ref{sec:6}--\ref{sec:7}).
\end{enumerate}

One important point about these results should be clear before quoting them. The learned figures in this paper are sensitive to revisions: between the two most recent runs, using the same seeds and data, the training-free results are identical, but the learned fold means can change by up to $0.172$---twice the $+0.085$ margin that the main comparison relies on (Section~\ref{sec:8.3}). The direction of this difference is stable in all our tests, but its size depends on the specific artifacts we provide, not just the method.

\paragraph{Relationship to the authors' prior work}

A previous conference paper \cite{yigit2025graph} introduced the preliminary multigraph formulation and deterministic quality model on synthetic topologies. This JSS manuscript substantially extends that work through incorporating the complete predictive HGT pathway with 16-dimensional QoS edge encoding (Section~\ref{sec:4}); inductive LOSO cross-validation (Section~\ref{sec:7.2}); zero-shot evaluation throughout five open-source systems (Section~\ref{sec:7.4}); empirical cost and sustainability characterization (Section~\ref{sec:7.5}); multi-oracle convergent validity and graph-view separation (Sections~\ref{sec:4.3}--\ref{sec:4.4}); and system-wide sensitivity analyses (Section~\ref{sec:7.3}). We retained formalisms from the conference paper only in restructured portions of Sections~\ref{sec:3} and~\ref{sec:5}. Two of its conclusions do not survive the re-measurement reported here. We state that explicitly rather than leaving the earlier record standing: the deterministic quality model is not a competitive ranking instrument on this corpus ($\rho = 0.205$, below unweighted centrality, Section~\ref{sec:7.1}), and its elicited weights rank worse than a uniform prior (Section~\ref{sec:7.3}). The corpus, the oracles, and the evaluation population all differ from \cite{yigit2025graph}, so these are corrections to what we now believe rather than a failure to reproduce a fixed experiment.

\subsection{Paper Organization}
\label{sec:1.6}

The rest of this paper is organized as follows: Section \ref{sec:2} reviews related work. Section \ref{sec:3} explains the SaG multigraph model and dependency projections. Section \ref{sec:4} describes the Heterogeneous Graph Transformer and simulation oracles. Section \ref{sec:5} presents the ISO/IEC-based explanation layer. Section \ref{sec:6} covers the experimental protocol, and Section \ref{sec:7} reports results for RQ1--RQ5. Section \ref{sec:8} discusses practical consequences, sustainability, threats to validity, and limitations. Section \ref{sec:9} concludes.
